% Transcription — fonds Alexandre Grothendieck, shelfmark 64, batch 7 (pages 121-140).
% Established from the facsimile archives/batches/64/batch-07.pdf (a copy of 64.pdf
% fetched through the project relay, no cover sheet: PDF page k = archivists' page 20(N-1)+k),
% per the skill transcribe-grothendieck. The folder is 156 pages in eight batches.
%
% Pass: Opus 5.5 (claude-opus-5-5), 2026-09-27 — first pass, unchecked against the pages by a human.
%
% Pages skipped: 121 (blank, show-through of 122), 123 (typed USTL thesis-defence notice),
%   126 (typed USTL seminar announcement), 129 (blank). Pages 127 and 128 are covers in his
%   hand; they get no \page{} and become section headings.
% His pagination: none visible in this batch (a circled, overwritten number top right of p. 140
%   is not read). Equations are numbered (1)-(28) from p. 130 on; (27) is used twice (p. 139).
% What the batch is about: pp. 122-125, a computation expressing a basis (f_1, f_2) in terms
%   of e_1-e_2 etc. for three distinct points xi_i (values of wp at 2-torsion points), with
%   Delta_0 = prod (xi_i - xi_j). Then « Approche transcendante »: complex structures on a real
%   plane V as non-real points of P(V)(C), elliptic curves over C as (Pi, omega, z), Torelli
%   rigidifications Z^2 -> Pi, the upper half-plane D, careful sign conventions (z = -(z_2/z_1)^-),
%   recovery of the usual action (az+b)/(cz+d), the theorem D -> M_{1,1}^an universal cover,
%   corollaries T_{1,1} = SL(2,Z), mapping class group of the pointed torus, Autext^+ of a group
%   of type (1,1); p. 140 starts on the tangent bundle along the unit section.
% Notation fixed once for the batch: \Pi for his doubly-barred lattice symbol (Pi = pi_1(E,0));
%   \mathfrak{D} for the half-plane; \mathbb{U} for the circle group; \mathcal{T} for his looped T
%   (Teichmüller-type group); \underline{\omega} for orientation sets as in batch 1.
% \ill{}: 72. \uncertain{}: 75.

\documentclass[11pt,a4paper]{article}
\input{../preamble/grothendieck}

\folder{64}
\batch{7}
\pages{121}{140}
\dating{[à partir de 1980- 1982]}
\watermark{Édition de démonstration}
\foldertitle{Modules courbes elliptiques en caractéristiques ≠2 : notes manuscrites (s.d.)}

\begin{document}

\page{122}
\[
  \begin{array}{ccc}
    \xi_1 & \xi_2 & \xi_3 \\
    \downarrow & \downarrow & \downarrow \\
    0 & 1 & \infty
  \end{array}
  \in k , \qquad \infty \mapsto t
\]
\note{feuillet pris en largeur. En tête, trois éléments $\xi_1, \xi_2, \xi_3$ de $k$ envoyés sur $0, 1, \infty$ ; une flèche courbe relie le $\infty$ écrit à droite à une lettre $t$ (lecture peu sûre). Au-dessous, un croquis : dans le plan, trois droites issues de l'origine passent par les points $\xi_1$, $\xi_3$, $\xi_2$ marqués sur une horizontale, soit les droites engendrées par les vecteurs ci-dessous.}
\[
  \begin{cases}
    (\xi_1, 1) \\
    (\xi_2, 1) \\
    (\xi_3, 1)
  \end{cases}
  \qquad
  \begin{array}{l}
    \lambda_1 \xi_1, \lambda_1 \\
    \lambda_2 \xi_2, \lambda_2 \\
    \lambda_3 \xi_3, \lambda_3
  \end{array}
\]
\[
  \left[\;
  \begin{aligned}
    &\textstyle\sum \lambda_i = 0 \\
    &\textstyle\sum \lambda_i \xi_i = 0
  \end{aligned}
  \right.
\]
$\xi_1, \xi_2, \xi_3$ \uncertain{deux à deux} distincts. Les \uncertain{déterminés} $(\lambda_i)$
\uncertain{sont déterminés} \add{mult.\ par} \uncertain{scalaire} \ill{}
\note{la ligne sous l'accolade est rapide ; le sens est que la solution $(\lambda_i)$ du système est déterminée à un multiple scalaire près.}

ou
\[
  \begin{aligned}
    \lambda_1 &= \xi_2 - \xi_3 \\
    \lambda_2 &= \xi_3 - \xi_1 \\
    \lambda_3 &= \xi_1 - \xi_2
  \end{aligned}
  \qquad
  \begin{cases}
    e_2 - e_3 = \bigl((\xi_2 - \xi_3)\xi_1,\ \xi_2 - \xi_3\bigr) \\
    e_3 - e_1 = \bigl((\xi_3 - \xi_1)\xi_2,\ \xi_3 - \xi_1\bigr) \\
    e_1 - e_2 = \bigl((\xi_1 - \xi_2)\xi_3,\ \xi_1 - \xi_2\bigr)
  \end{cases}
\]
\note{entre les deux colonnes, une colonne $\xi_1, \xi_2, \xi_3$ séparée par un trait vertical. Au bas de la page, un fragment biffé : \struck{$(\xi_1, 1) \mapsto$} \ldots\ \struck{$f_1 \to$}, puis « $\xi_1 f_1 +$ », inachevé.}

\page{124}
\[
  \begin{aligned}
    \xi_1 f_1 + f_2 &= \frac{1}{\xi_2 - \xi_3}\,(e_2 - e_3) \\
    \xi_2 f_1 + f_2 &= \frac{1}{\xi_3 - \xi_1}\,(e_3 - e_1) \\
    \xi_3 f_1 + f_2 &= \frac{1}{\xi_1 - \xi_2}\,(e_1 - e_2)
  \end{aligned}
\]
\[
  (\xi_1 - \xi_2) f_1 = \frac{1}{\xi_3 - \xi_1}\, e_1 + \frac{1}{\xi_2 - \xi_3}\, e_2
  - \Bigl(\frac{1}{\xi_3 - \xi_1} + \frac{1}{\xi_2 - \xi_3}\Bigr) e_3
\]
\[
  f_1 = \frac{1}{(\xi_1 - \xi_2)(\xi_3 - \xi_1)}\, e_1
  + \frac{1}{(\xi_1 - \xi_2)(\xi_2 - \xi_3)}\, e_2
  - \frac{1}{\xi_1 - \xi_2}\,(\;\cdot\;)\, e_3
\]
\note{la parenthèse devant $e_3$ reprend, par une accolade tracée au-dessus, la somme de la ligne précédente ; son contenu n'est pas écrit.}
\[
  \begin{aligned}
    f_2 &= \frac{1}{\xi_2 - \xi_3}\,(e_2 - e_3) - \xi_1 f_1 \\
    &= -\frac{\xi_1}{(\xi_1 - \xi_2)(\xi_3 - \xi_1)}\, e_1
      + \Bigl(\frac{1}{\xi_2 - \xi_3} - \frac{\xi_1}{(\xi_1 - \xi_2)(\xi_2 - \xi_3)}\Bigr) e_2 \\
    &\qquad + \Bigl[-\frac{1}{\xi_2 - \xi_3}
      + \frac{\xi_1}{\xi_1 - \xi_2}\Bigl(\frac{1}{\xi_3 - \xi_1} + \frac{1}{\xi_2 - \xi_3}\Bigr)\Bigr] e_3
  \end{aligned}
\]
\[
  \begin{aligned}
    (\xi_1 - \xi_2)(\xi_2 - \xi_3)(\xi_3 - \xi_1)\, f_2
    &= -\xi_1(\xi_2 - \xi_3)\, e_1
      + \bigl[(\xi_1 - \xi_2)(\xi_3 - \xi_1) - \xi_1(\xi_3 - \xi_1)\bigr] e_2 \\
    &\qquad + \bigl[-(\xi_1 - \xi_2)(\xi_3 - \xi_1) + \xi_1(\xi_2 - \xi_3) + \xi_1(\xi_3 - \xi_1)\bigr] e_3 \\
    &= -\xi_1(\xi_2 - \xi_3)\, e_1 - \xi_2(\xi_3 - \xi_1)\, e_2 - \xi_3(\xi_1 - \xi_2)\, e_3
  \end{aligned}
\]
\note{le signe devant le crochet du coefficient de $e_3$ est surchargé d'une tache ; la lecture « $-$ » est celle que donne le calcul. Des points et des traits obliques pointent les termes qui se simplifient ; les deux derniers signes « $-$ » de la dernière ligne sont repassés.}

\page{125}
\struck{$\xi_1(\xi_2 - \xi_3) \quad \xi_2(\xi_3 - \xi_1)$}
\[
  (\xi_1 - \xi_2)(\xi_2 - \xi_3)(\xi_3 - \xi_1) = \Delta_0
\]
\[
  \boxed{\;
  \begin{aligned}
    -\Delta_0 f_2 &= \xi_1(\xi_2 - \xi_3)\, e_1 + \xi_2(\xi_3 - \xi_1)\, e_2 + \xi_3(\xi_1 - \xi_2)\, e_3 \\
    \Delta_0 f_1 &= (\xi_2 - \xi_3)\, e_1 + (\xi_3 - \xi_1)\, e_2 + (\xi_1 - \xi_2)\, e_3
  \end{aligned}\;}
\]
$t$ correspond donc : $\sum (\xi_j - \xi_k)\, e_i$
\[
  = \sum_i \bigl(\wp(\varepsilon_j) - \wp(\varepsilon_k)\bigr)\, e_i
\]
\note{feuillet pris en largeur. Les indices $(i, j, k)$ parcourent les permutations circulaires de $(1, 2, 3)$ ; les $\varepsilon$ sont lus d'après le contexte (valeurs de $\wp$ aux points d'ordre $2$), le premier est surchargé.}

\section{Relations entre $M_{11}$ et $M_{04}$}
\note{titre de l'auteur, seul sur un feuillet par ailleurs blanc (p.~127). Aucune page de ce lot ne le suit avant la couverture suivante.}

\section{Approche transcendante}
\note{titre de l'auteur, seul sur un feuillet par ailleurs blanc (p.~128), précédé de la mention « (21 p) » entourée ; p.~129 est blanche.}

\page{130}
\subsection{Modules des courbes elliptiques : théorie transcendante}
\note{titre de l'auteur, en tête de page.}

La donnée d'une courbe elliptique \add{$E$} sur $\mathbb{C}$ revient à celle d'un
\add{$\mathbb{C}$-}espace vectoriel de rang $1$, savoir $t_E$
\[
  V = t_E \qquad \text{(espace tangent à l'origine)} \tag{1}
\]
et d'un « réseau »
\[
  \Pi \subset V \tag{2}
\]
(sous-\struck{groupe additif} \add{$\mathbb{Z}$-module, i.e.} \struck{discret}
\add{discret, libre} de rang $2$). La donnée du \struck{groupe} $\mathbb{Z}$-module
$\Pi$ permet de récupérer $V$ en tant qu'espace vectoriel réel \add{(de rang $2$)}, par
\[
  V \simeq \Pi \otimes_{\mathbb{Z}} \mathbb{R} \overset{\text{déf}}{=} \Pi_{\mathbb{R}} \tag{3}
\]
et $E$ en tant que groupe de Lie réel, par
\[
  E \simeq V/\Pi \simeq \Pi_{\mathbb{R}}/\Pi , \tag{4}
\]
tandis que la structure complexe de $E$ est déterminée par la donnée d'une structure
d'espace vectoriel complexe sur $V$ par passage aux quotients \struck{\uncertain{de groupes}}.
\note{le réseau est noté par un symbole à deux jambages barrés, qui ressemble à un $H$ ou à un $\Pi$ doublement barré ; le texte l'appelle « le $\mathbb{Z}$-module $\Pi$ », et on le transcrit $\Pi$ partout dans ce lot.}

Étant donné un espace vectoriel réel $V$, la donnée d'une structure complexe sur $V$
(i.e.\ d'un automorphisme $u$ de $V$ de carré $-\mathrm{id}_V$) équivaut à ceci : celle
d'un sous-espace vectoriel \add{$\mathbb{C}$-}$L$ de \struck{\ill{}} 
\[
  V_{\mathbb{C}} = V \otimes_{\mathbb{R}} \mathbb{C} \tag{5}
\]
tel que la projection canonique
\[
  \begin{aligned}
    V_{\mathbb{C}} &\longrightarrow V = V_{\mathbb{R}} \\
    z &\longmapsto \tfrac{1}{2}(z + \bar z) = \Re z
  \end{aligned} \tag{6}
\]
\struck{soit} \add{induise} un isomorphisme d'espaces vectoriels réels
\struck{Pour un}
\[
  \varphi : L \xrightarrow{\;\sim\;} V \qquad (\mathbb{R}\text{-iso}) \tag{7}
\]

\page{131}
\add{(}à un tel $L$, on associe la structure complexe de $V$ \struck{\ill{}, \ill{},
conjuguée de celle déduite de \ill{}} \struck{\ill{} de celle de \ill{} par transport
de} \add{$V$} qui est conjuguée de celle déduite de la str.\ complexe de $L$ par
transport de structure via (7), i.e.\ on pose
\[
  u x = \varphi(-i\,\varphi^{-1}(x)) = -\varphi(i\,\varphi(x)) \tag{8}
\]
$u$ étant la multiplication par $i$ pour la structure complexe envisagée sur $V$.
\marginal{\uncertain{cette convention} \ill{} \ill{} par \uncertain{les} \ill{} \ill{}
\ill{} \uncertain{courbes elliptiques} \ill{} \uncertain{de fait qui} \ill{} \ill{}
\ill{} \uncertain{introduit} \ill{} \ill{} \ill{}}
\note{la note marginale est écrite de biais dans l'angle supérieur gauche, sur plusieurs lignes, et reliée par un trait au début de la page ; elle justifie le choix de la structure \emph{conjuguée}, mais n'est lue que par bribes. Le signe « $-$ » dans $\varphi(-i\,\varphi^{-1}(x))$ est surchargé ; la seconde égalité, avec $\varphi(x)$ au lieu de $\varphi^{-1}(x)$, est telle qu'écrite.}

Inversement, si $V$ est muni d'une structure complexe, \struck{l'applic.} faisant de
$V$ un espace vectoriel complexe $V_c$, alors l'application \add{identique}
$\mathbb{R}$-linéaire $V \to V_c$ se prolonge, par la propriété universelle de
$V_{\mathbb{C}}$, en une application $\mathbb{C}$-linéaire
\[
  \psi : V_{\mathbb{C}} \longrightarrow V_c \tag{9}
\]
de sorte que
\[
  \begin{cases}
    \psi(i x) = u x \\
    \psi(x) = x
  \end{cases}
  \qquad \forall\, x \in V = V_{\mathbb{R}}, \tag{10}
\]
donc
\[
  \psi(x + i y) = x + u y . \tag{11}
\]
\struck{Considérons} \struck{Prenons} Considérons
\[
  \begin{aligned}
    L' = \operatorname{Ker} \psi
    &= \{\, x + i y \mid x, y \in V_{\mathbb{R}},\ x + u y = 0 \ \text{ i.e. } y = u x \,\} \\
    &= \{\, x + i u(x) \mid x \in V_{\mathbb{R}} \,\}
  \end{aligned} \tag{12}
\]
et considérons l'application correspondante (6)
\[
  \varphi : \underset{\substack{\| \\ x + i u(x)}}{z} \longmapsto \tfrac{1}{2}(z + \bar z) = \Re z = x : L \longrightarrow V .
\]
\note{dans (12), la lettre primée est surchargée ($L$ puis $L'$) ; « $y = ux$ » récrit sur un premier essai. Pour $x + uy = 0$ on a en fait $y = ux$ puisque $u^2 = -1$ ; la ligne d'accolade du second membre est surchargée d'un premier essai biffé.}

On voit alors que cette application est bijective, et \uncertain{antilinéaire},
\struck{\ill{} la \uncertain{conjuguée} \ill{} de la structure complexe \ill{} pour $V$
\ill{} \ill{} $L = L' = \operatorname{Ker}\psi \subset V_{\mathbb{C}}$, \ill{} $\mathbb{R}$}
\marginal{\uncertain{mais pas} \ill{}, \uncertain{alors} $V \to V_c$ \uncertain{antilinéaire}}
\note{deux lignes et demie sont biffées ici ; on n'en lit que des fragments. La note marginale gauche, reliée à la ligne suivante, précise le caractère antilinéaire de l'application.}

(13) [On peut dire aussi, que la donnée d'une structure complexe sur $V$ équivaut à :

1°) La donnée d'un \struck{\ill{}} sous-groupe compact \struck{maximal} $U$ \struck{de}
$\simeq \mathbb{C}^{1}$ \uncertain{de} $\operatorname{Aut}_{\mathbb{R}}(V)$, \struck{i.e.\ d'une}
\note{« $\simeq \mathbb{C}^{1}$ » est une lecture peu sûre d'un groupe surchargé écrit au-dessus de la ligne.}

\page{132}
\struck{\ill{} \ill{}} \add{\uncertain{isomorphe}} \uncertain{isomorphe} : $\mathbb{U} = \{\, z \in \mathbb{C} \mid |z| = 1 \,\}$,
tel que l'action sur $V_{\mathbb{C}}$ ait exactement $2$ sous-espaces propres \add{\uncertain{propres}}
\marginal{(\uncertain{i.e.}\ \ill{} \uncertain{dominantes} \ill{} si $\operatorname{rg}_{\mathbb{R}} V = 2$)}

2°) d'un des deux éléments d'ordre $4$ de $\mathbb{U}$, ce qui revient au \ill{} d'une
orientation de $\mathbb{U}$, ou encore \add{(si $\operatorname{rg}_{\mathbb{R}} V = 2$, plus
généralement si $\operatorname{rg}_{\mathbb{R}} V = 2n$, $n$ impair)} d'une orientation de $V$
--- l'élément d'ordre $4$ de $\mathbb{U}$ associé à une orientation $\omega$ de
\add{$\mathbb{U}$ (ou de} $V$) sera noté $u_\omega$, c'est donc la rotation de $+\pi/2$ pour
l'orientation \add{$\omega$} de $V$ et la structure \struck{complexe} \add{\uncertain{symplectique}}
sur $V$ choisie invariante sous $\mathbb{U}$ (\underline{\uncertain{cas}} \add{si
$\operatorname{rg}_{\mathbb{R}} V = 2$,} une telle \struck{structure} \add{symplectique} est
unique : scalaire multiplicatif près).
\marginal{\uncertain{on engendre} \ill{} \uncertain{d'un isomorphisme} $\mathbb{U} \xrightarrow{\sim} \mathbb{U}$, $\mathbb{U} \xrightarrow{\sim} \mathbb{U}$ \ill{}}
\note{le groupe est écrit avec une lettre $U$ à double jambage ; on le note $\mathbb{U}$. Les deux notes marginales gauches, écrites de biais, ne sont lues qu'en partie.}

Alors $V_{\mathbb{C}}$ sous l'action de $\mathbb{U}$ se décompose en deux sous-espaces propres
\[
  V_{\mathbb{C}} = \bigoplus_{\omega \in \underline{\omega}} L_\omega \tag{14}
\]
où pour $\omega \in \underline{\omega}(\mathbb{U}) = \underline{\omega}$, $L_\omega$ est caractérisé par
le fait que l'isomorphisme
\[
  \varphi_\omega : \mathbb{U} \xrightarrow{\;\sim\;} \mathbb{U} \qquad
  g(z) = \varphi_\omega(g)\cdot z \ \text{ pour } z \in L_\omega \tag{15}
\]
défini par l'action de $\mathbb{U}$ sur $L_\omega$, correspond à l'orientation $\omega$.
La donnée d'une structure complexe sur $V$ \add{compatible avec $\mathbb{U}$} équivaut alors à celle
d'un des deux isomorphismes $\varphi_\omega$ permettant de faire opérer $\mathbb{U}$ sur $V$,
\uncertain{donc aussi} $\mathbb{R}^{\times}_{+} \times \mathbb{U} \simeq \mathbb{C}^{\times}$,
\struck{\ill{}} $\mathbb{C}$ \ldots). Ceci dit, pour \struck{\ill{}} $z \in L_\omega$, on aura
\[
  \Re(\varphi_\omega(g) z) = \Re(g z) = g\, \Re(z) = \text{produit de } \Re(z) \text{ par } \varphi_\omega(g)
  \text{ pour la structure complexe de } V \text{ associée à } \varphi_\omega
\]
\note{le numéro de (14) est surchargé (13 corrigé en 14). La parenthèse fermante après « $\mathbb{C} \ldots$ » répond à une parenthèse ouverte que la page ne montre pas.}

\page{133}
i.e.
\[
  \begin{aligned}
    L_\omega &\xrightarrow{\;\sim\;} V \\
    z &\longmapsto \Re z
  \end{aligned}
\]
est un isomorphisme d'espaces vectoriels complexes, quand on munit $V$ de la structure
complexe $c_\omega$.~]
\marginal{Mais il y a lieu d'associer à cette structure \uncertain{complexe} sur $V$ \ill{} \uncertain{non} par $L_\omega$, mais $L_{\omega'}$ \uncertain{plutôt} \ldots}
\note{le crochet fermant clôt celui ouvert en (13), p.~131. La note marginale, de biais dans la marge gauche, n'est lue qu'en partie.}

Dans le cas qui nous occupe, $\operatorname{rg}_{\mathbb{R}} V = 2$, les \struck{droites}
sous-espaces vectoriels envisagés de $V_{\mathbb{C}}$ sont de \struck{\ill{}} rang $1$,
\add{\uncertain{cotangents} $1$,} et peuvent s'interpréter comme les points de
\[
  \Sigma = \Sigma_V = P(\mathbb{C}), \qquad \text{où } P = \mathbb{P}(V) \tag{16}
\]
est la \uncertain{relative} droite projective réelle associée à $V$. Pour que
\[
  \begin{aligned}
    L &\longrightarrow V \\
    z &\longmapsto \Re z
  \end{aligned}
\]
soit un isomorphisme, \struck{i.e.} pour que $L$ corresponde à une structure complexe
\struck{de} \uncertain{sur} $V$, \add{$L \cap V = \{0\}$ i.e.\ $L$} il faut et il suffit que
\struck{le point $z$ correspondant de $\Sigma$ soit \uncertain{non réel}} $L$ n'est pas
complexifié d'une droite $L_0 \subset V_{\mathbb{R}}$, i.e.\ que le point $z \in \Sigma$ qui
correspond à $L$ \add{ne} soit pas dans \struck{$\Sigma(\mathbb{R})$}
\[
  \Sigma_{\mathbb{R}} = P(\mathbb{R})
\]
ou encore qu'on ait
\[
  z \in \Sigma \setminus \Sigma_{\mathbb{R}} . \tag{17}
\]
\struck{En résumé} Ainsi, les structures complexes
\note{« $V_{\mathbb{R}}$ » : l'indice de $V$ est surchargé. La phrase continue p.~134.}

\page{134}
sur $V$ sont en correspondance biunivoque avec l'ens.\ $\Sigma \setminus \Sigma_{\mathbb{R}}$
des points non réels de $\Sigma$.

Notons que l'ens.\ $\Sigma \setminus \Sigma_{\mathbb{R}}$ a deux composantes connexes
\add{$\Sigma^*_\omega$}, qui correspondent (par Stokes) aux deux orientations de
$\Sigma_{\mathbb{R}}$, ou encore aux deux orientations de $V$. La composante connexe d'un
$z \in \Sigma \setminus \Sigma_{\mathbb{R}}$ correspond à l'orientation de $V$ définie de
façon évidente par la structure complexe envisagée sur $V$ (définie par $z$).
\marginal{\uncertain{Il faut préciser} quelle \uncertain{convention} \ill{} \uncertain{vers le} bas}
\note{la note marginale, écrite verticalement et reliée par un trait au paragraphe, n'est lue qu'en partie.}

Revenant à la description transcendante des courbes elliptiques, on voit que
\struck{une telle} \add{la donnée d'une telle} courbe $E$ revient à la donnée
\begin{enumerate}
  \item[a)] D'un $\mathbb{Z}$-module libre de rang $2$ $\Pi$
  \item[b)] D'une \struck{droite} orientation $\omega$ de $V = \Pi \otimes_{\mathbb{Z}} \mathbb{R}$,
    ou ce qui revient au même, d'une des deux bases de $\bigwedge^2_{\mathbb{Z}} \Pi$ (i.e.\ d'une
    restriction du groupe structural $\mathrm{GL}(2, \mathbb{Z})$ de $\Pi$ : $\mathrm{SL}(2, \mathbb{Z})$),
    \struck{soit} notée aussi $\omega$
  \item[c)] D'un point $z \in \Sigma_\omega(V)$,
\end{enumerate}
\marginal{On a $\Pi = \pi_1(E, 0_E)$, $V = \Pi \otimes_{\mathbb{Z}} \mathbb{R} \simeq t_E$}

Une \underline{rigidification de Torelli} de la courbe elliptique $E$ est la donnée d'un
isomorphisme
\note{« $\Sigma_\omega(V)$ » : le $\Sigma$ est surchargé ; il s'agit de la composante $\Sigma^*_\omega$ introduite plus haut. La phrase continue p.~135.}

\page{135}
\[
  \mathbb{Z}^2 \longrightarrow \Pi \tag{18}
\]
compatible avec les orientations, i.e.\ tel que l'on ait
\[
  z_1 \wedge z_2 = \omega \tag{19}
\]
où $z_1, z_2$ sont les images par (18) des éléments de la base canonique de $\mathbb{Z}^2$.
Cet isomorphisme (18) définit un isomorphisme correspondant
\[
  \mathbb{R}^2 \xrightarrow{\;\sim\;} V \tag{20}
\]
d'où des isomorphismes
\[
  \begin{cases}
    \mathbb{P}^1_{\mathbb{R}} \simeq \mathbb{P}^1(V), \quad \mathbb{P}^1(\mathbb{C}) \simeq \Sigma_V \\
    \mathfrak{D} \xrightarrow{\;\sim\;} \Sigma^*_{V,\omega}
  \end{cases} \tag{21}
\]
où $\mathfrak{D}$ est le demi-plan de Poincaré, i.e.\ \struck{\ill{}} la variété complexe des
$z \in \mathbb{C}$ tels que \struck{$\Im z$} $\Im z > 0$
\[
  \mathfrak{D} = \{\, z \in \mathbb{C} \mid \Im z > 0 \,\}, \tag{22}
\]
identifié à $\Sigma^*_{\mathbb{R}^2, \omega_0}$ en identifiant $z \in \mathfrak{D}$ à la droite
\[
  L_z = \mathbb{C}\cdot\underset{z e_1 + e_2}{\underbrace{(z, 1)}} . \tag{23}
\]
\marginal{si $(e_1, e_2)$ base canonique de $\mathbb{R}^2$ \ill{} : \uncertain{associant} à $z \in \mathbb{C}$ la droite $\mathbb{C}(1, z)$, alors \uncertain{correspond} \ill{} $\Sigma^*_{\mathbb{R}^2, \omega_0}$ \uncertain{donnerait} $\mathfrak{D} = \mathfrak{D}^+$ \ill{}}
\note{l'indice de $L_z$ dans (23) est surchargé. La note marginale, écrite de biais dans l'angle inférieur gauche, envisage l'autre convention $\mathbb{C}(1, z)$ ; elle n'est lue qu'en partie.}

Sur l'ens.\ des \struck{\ill{}} \add{classes d'isomorphie de} courbes elliptiques
$E = (\Pi, \omega, L \subset \Pi \otimes_{\mathbb{Z}} \mathbb{C})$, avec rigidification de Torelli
$\mathbb{Z}^2 \simeq \Pi$, le groupe $\mathrm{SL}(2, \mathbb{Z})$ opère : \struck{droite} \add{à droite}
de façon évidente, par composition \add{à droite} de (18) avec les automorphismes de $\mathbb{Z}^2$.
\struck{\uncertain{Quant}} \uncertain{On va}
\marginal{Ainsi, les classes d'isomorphie de courbes elliptiques avec rigidification de Torelli sont en correspondance biunivoque avec \uncertain{l'ensemble} \uncertain{des pts} de $\mathfrak{D}$.}
\note{la note de droite, insérée entre (23) et la phrase suivante, est de l'auteur ; son point d'insertion n'est pas marqué.}

\page{136}
expliciter cette action en tant qu'action sur $\mathfrak{D}$ (défini par (22)).

Si une structure complexe sur $\mathbb{R}^2$ est définie par l'élément $z \in \mathfrak{D}$,
$z = x + iy$, alors l'application \add{$\mathbb{C}$-linéaire}
$\varphi : L_z = \mathbb{C}\cdot(z, 1) \longrightarrow V$ s'identifie à :
\[
  \underset{\substack{\| \\ a + ib}}{\lambda} \longmapsto
  \Re\bigl(\underbrace{\lambda(z e_1 + e_2)}_{[(ax - by) + i(ay + bx)] e_1 + (a + ib) e_2}\bigr)
  = (ax - by) e_1 + a e_2
\]
les images inverses $\lambda_1, \lambda_2$ de $e_1, e_2$ par cette \uncertain{app.}\ sont données par
\[
  \lambda_1 = -\frac{i}{y} \qquad \lambda_2 = 1 + i\,\frac{x}{y}
\]
d'où
\[
  \lambda_2 / \lambda_1 = i(y + ix) = -x + iy
\]
donc on a
\[
  \underset{\substack{\text{quotient dans } \mathbb{C}^* \\ \text{au sens de la} \\
  \text{structure compl.\ de } \mathbb{R}^2 \\ \text{définie par } z = x + iy \in \mathfrak{D}}}
  {\underbrace{e_2 / e_1}}
  = (-x + iy)^{-} = -x - iy = -z \ \in \mathfrak{D}^{-}
\]
\note{dans la définition de $\varphi$, un premier argument de $\Re$ est biffé et surchargé. Le « $^{-}$ » après $(-x+iy)$ est la barre de conjugaison. L'auteur écrit $\lambda_2/\lambda_1 = i(y+ix)$ ; le calcul donne en fait $(y+ix)/(-i) = i(y+ix)$, ce qui est cohérent.}

Donc si une structure de Torelli sur $E$ est définie par la \struck{\uncertain{donnée}}
\add{base} $(z_1, z_2)$ de $\Pi \subset t_E$, alors l'élément $z \in \mathfrak{D}$ correspondant est
donné par
\struck{$z_2/z_1 = -x + iy = -z$}
\[
  z = -(z_2 / z_1) \tag{24}
\]
\struck{où $z = x + iy$}.
\note{la ligne (24) est en partie biffée : l'auteur a d'abord écrit $z_2/z_1 = -x+iy = -z$, puis a récrit au-dessus $z = -(z_2/z_1)$ ; « où $z = x+iy$ » est biffé.}

Soit alors
\[
  g \in \mathrm{SL}(2, \mathbb{Z}) \qquad g = \begin{pmatrix} a & b \\ c & d \end{pmatrix}
  \qquad a, b, c, d \in \mathbb{Z},\ ad - bc = 1
\]
et soit \struck{\ill{}}
\[
  u : \mathbb{Z}^2 \xrightarrow{\;\sim\;} \Pi
\]
une structure de Torelli, correspondant à :

\page{137}
\[
  z_1 = u(e_1), \quad z_2 = u(e_2) \in \Pi \subset V
\]
alors la structure de Torelli $u \circ g$ correspond à $z'_1, z'_2 \in \Pi$ donnés par
\[
  \begin{aligned}
    z'_1 &= u\,g(e_1) = u(a e_1 + c e_2) = a z_1 + c z_2 \\
    z'_2 &= u\,g(e_2) = u(b e_1 + d e_2) = b z_1 + d z_2 ,
  \end{aligned}
\]
qui correspond donc à un élément
\[
  z' = x' + i y' \in \mathfrak{D}'
\]
donné par
\[
  z' = -\frac{z'_2}{z'_1} = -\frac{b z_1 + d z_2}{a z_1 + c z_2}
  = -\frac{d \frac{z_2}{z_1} + b}{c \frac{z_2}{z_1} + a}
  = \frac{d z - b}{-c z + a} \tag{25}
\]
où \struck{\ill{}}. Si on \struck{\ill{}} préfère expliciter une action à gauche
$u \mapsto u \circ g^{-1}$ de $\mathrm{SL}(2, \mathbb{Z})$ sur l'ens.\ des structures de Torelli
\uncertain{mod.\ iso}, on note que
\[
  g^{-1} = \begin{pmatrix} d & -b \\ -c & a \end{pmatrix}
\]
\struck{d'où} d'où immédiatement, pour $z' = x' + i y'$ correspondant à $u \circ g^{-1}$ :
\[
  z' = \frac{a z + b}{c z + d} \tag{25 bis}
\]
\underline{victoire}, c'est la formule habituelle,
\note{à côté de (25), un premier essai surchargé « $z' = \ldots$ » est biffé, avec au-dessous « $z' = -\bar z'$ » partiellement biffé ; le signe de la dernière fraction de (25) et la lettre $z$ au numérateur sont surchargés. Avant (25 bis), une ligne entière est biffée : $\tilde z' = -\bar z' = \frac{a\tilde z - b}{-c\tilde z + d} = \frac{-a\bar z - b}{-c\bar z - d} = \frac{a\bar z + b}{c\bar z + d}$, avec un numéro « (25 bis) » biffé à gauche. Dans la marge gauche, un schéma de calcul croisé (les produits $-dz$, $-c$, $b$, $a$, $x$, reliés par des traits obliques) aboutit à un « (25) » entouré ; il n'est pas transcrit plus avant.}

mais on fera attention que le nb complexe $z \in \mathfrak{D}$ attaché à la structure de Torelli
définie par \struck{$z_1, z_2 \in t_E$} \add{les deux} éléments $z_1, z_2$ de la droite complexe
\struck{\ill{}} $= t_E$ est donnée, non
\note{la phrase continue p.~138.}

\page{138}
\struck{par la formule $z_2/z_1$, qui \uncertain{naïvement}} \uncertain{paraîtrait} \uncertain{naturelle}
\[
  z = z_2 / z_1 ,
\]
mais par (24) \struck{\ill{}} i.e.\ par
\[
  z = -(z_2 / z_1)^{-} = -\bar z_2 / \bar z_1 \in \mathfrak{D}
\]
(\struck{\ill{}} si on voulait \add{\uncertain{associer}} répérer la structure de Torelli \uncertain{sur $\mathbb{R}^2$}
par le nb $z \in \mathfrak{D}$ choisi, mais \uncertain{par} le nombre complexe $\tau = z_2 / z_1$
lui-même, \uncertain{il faudrait} \ill{} $z$ \uncertain{par} à l'axe imaginaire), \ill{} \uncertain{donnerait}
une formule de transformation d'opération de $\mathrm{SL}(2, \mathbb{Z})$ sur $\mathfrak{D}$ \uncertain{distincte}
de la formule habituelle.
\marginal{Donc cela montre que la condition $z \in \mathfrak{D}^+$ \uncertain{signifie que} $(z_1, z_2)$ \uncertain{est une base} \uncertain{directe} \ill{} de $t_E$ \ill{} \ill{} \uncertain{Torelli} \ill{} \ill{} \ill{} \uncertain{complexe} \ill{} ; \uncertain{inverse} \ill{} \uncertain{voulait} faire \ill{} $z \in \mathfrak{D}^+$ \ill{} \ill{} (23) \ill{} $\mathbb{C}\cdot(1, z)$ \ill{} $L_z$ \ill{} $\Sigma^*_{\mathbb{R}^2, \omega_0}$, \ill{} $z = -(z_2/z_1)^{-}$, \uncertain{il faudrait} \ill{} \uncertain{naturelles} \ldots}
\note{la marge gauche porte, sur toute la moitié supérieure de la page, une longue note écrite de biais, en partie biffée, qui reprend la question des conventions (base directe, choix de $\mathbb{C}(1, z)$ au lieu de $\mathbb{C}(z, 1)$ en (23)) ; on n'en donne que des bribes. Plusieurs mots du corps, entre les lignes, sont aussi incertains ; le sens général est que la formule naïve $\tau = z_2/z_1$ conduirait à une action de $\mathrm{SL}(2, \mathbb{Z})$ différente de la formule habituelle.}

\underline{Théorème}. \struck{On trouve ici la} $\mathfrak{D}$ \add{$= \{\, z \in \mathbb{C} \mid \Im z > 0 \,\}$,}
muni de la famille \add{\uncertain{analytique}} de courbes elliptiques
\[
  E_z = \mathbb{C} / (\mathbb{Z}\cdot 1 + \mathbb{Z}\cdot(-\bar z))
\]
est un revêtement universel de $M^{\mathrm{an}}_{1,1}$, correspondant à la classification des courbes
elliptiques avec rigidification de Torelli ; le groupe des $M^{\mathrm{an}}_{1,1}$-automorphismes de
$\mathfrak{D}$ (opérant à droite) s'identifie au groupe $\mathrm{SL}(2, \mathbb{Z})$, opérant sur
$\mathfrak{D}$ par
\[
  g z = \frac{a z + b}{c z + d} \qquad \text{pour } g = \begin{pmatrix} a & b \\ c & d \end{pmatrix} \in \mathrm{SL}(2, \mathbb{Z}),
\]
qui correspond aussi à l'opération à gauche \uncertain{naturelle}
\note{l'exposant de $M_{1,1}$, deux points surmontés d'un trait, est lu « an ». Une parenthèse ouvrante isolée précède $E_z$. La phrase continue p.~139.}

\page{139}
de $\mathrm{SL}(2, \mathbb{Z})$ sur les structures de Torelli, par
\[
  g \cdot u = u \circ g^{-1} \qquad (u : \mathbb{Z}^2 \xrightarrow{\;\sim\;} \Pi_E).
\]

\underline{Corollaire 1}. On a un iso.\ canonique
\[
  \mathcal{T}_{1,1} \simeq \mathrm{SL}(2, \mathbb{Z}) \tag{26}
\]
(où $\mathcal{T}_{1,1} = \pi_1(M_{1,1})$ est calculé relativement au revêtement universel choisi
$\mathfrak{D}$ de $M_{1,1}$). On a un iso.\ \add{can.}
\[
  M_{1,1} \simeq (\mathfrak{D}, \mathrm{SL}(2, \mathbb{Z})) \tag{27}
\]
\note{le groupe noté $\mathcal{T}$ est écrit comme un T cursif bouclé ; nous le transcrivons $\mathcal{T}$. Dans (27), la lettre $\mathfrak{D}$ est surchargée. Le numéro (27) est repris plus bas pour une autre formule.}

Utilisant les grands théorèmes \uncertain{orthodoxes}-\uncertain{piles}, on en déduit

\underline{Corollaire 2}. Soit \struck{$(E)$} $(E, a)$ une surface \add{\uncertain{pointée}} orientée
compacte connexe de genre $1$, et soit $\mathcal{T}$ le groupe des classes d'isotopie des
automorphismes de $(E, a)$. \struck{Alors} \struck{Considérons} \add{Alors} l'homomorphisme naturel
\[
  \mathcal{T} \longrightarrow \operatorname{Aut}^+_{\mathbb{Z}}(\Pi) \qquad \text{où } \Pi = \pi_1(E, a) \tag{27}
\]
est un isomorphisme. \struck{\ill{}}

\underline{Corollaire 3}. Soit $\pi$ un groupe \uncertain{discret} de type $(1, 1)$, $H = \pi_{\mathrm{ab}}$ (qui est un
$\mathbb{Z}$-module libre de rang $1$), $\mathcal{T} = \operatorname{Autext}^+_{\mathrm{loc}}(\pi)$,
\struck{considérons} alors l'homom.
\[
  \mathcal{T} \longrightarrow \operatorname{Aut}^+_{\mathbb{Z}}(H) \tag{28}
\]
est iso.
\marginal{Il faudrait expliciter \uncertain{l'action} de $\operatorname{Aut}^+_{\mathbb{Z}}(H)$ \ill{} \ill{} \uncertain{extérieure} sur $\pi$ \ill{} $\operatorname{Aut}^+_{\mathbb{Z}}(H)$ \ldots}
\note{le mot lu « discret » est court et incertain ; l'auteur écrit « rang 1 » pour $H$ ; l'abélianisé d'un groupe de type $(1,1)$ est de rang $2$, et le « 1 » est sans doute un lapsus, laissé tel quel. Dans (27) de cette page, le groupe est noté par le même symbole doublement barré que le réseau $\Pi$. La note marginale, de biais dans l'angle inférieur gauche et reliée par un trait au Corollaire 3, n'est lue qu'en partie.}

\page{140}
Nous pouvons expliciter le fibré tangent le long des \struck{fibres} \add{la section unité} de la
courbe elliptique universelle rigidifiée à la Torelli, comme étant le fibré formé des $t_E$, qui
du point de vue fibré vectoriel réel s'identifie au \struck{fibré} $\mathrm{SL}(2, \mathbb{Z})$-fibré
sur $\mathfrak{D}$ défini par l'opération standard de $\mathrm{SL}(2, \mathbb{Z})$ sur $\mathbb{R}^2$.
Le fibré principal associé \add{$P\mathfrak{D}$,} s'obtient (du point de vue de la seule structure
topologique) en remplaçant \uncertain{dans ce} \uncertain{qui précède} $\mathbb{R}^2$ par
$(\mathbb{R}^2)^*$. Pour en trouver un revêtement universel, il suffit de choisir un revêtement
universel de $(\mathbb{R}^2)^* \simeq \mathbb{C}^*$, p.\ ex.\ le revêtement par l'exponentielle
complexe
\[
  z \longmapsto \exp(2 i \pi z) : \mathbb{C} \longrightarrow \mathbb{C}^*
\]
\note{en haut à droite, un numéro entouré et surchargé, non lu. Le reste de la page est blanc.}

\end{document}
